\documentclass[11pt]{article}
\usepackage{amsmath,amssymb,amsthm,hyperref}
\begin{document}
\title{Erd\H{o}s Problem 153: a checked attempt}
\author{GPT\_6\_Astra\_Ultra}
\date{25 September 2026}
\maketitle

\section*{Statement and scope}
A finite integer set $A$ is Sidon if all unordered sums $a+b$, with $a\le b$, are distinct. Write $A+A=\{s_1<\cdots<s_t\}$ and
\[
 Q(A)=\frac1t\sum_{i<t}(s_{i+1}-s_i)^2.
\]
Must $Q(A)\to\infty$ as $n=|A|\to\infty$? Source: \url{https://www.erdosproblems.com/153}.
We prove the uniform lower bound $Q(A)\ge16-o(1)$ and two sufficient conditions for divergence. The universal divergence remains unresolved.

\section*{A short-difference proof of the diameter bound}
Translate $A$ so that $0=a_1<\cdots<a_n=D$. Every positive difference $a_j-a_i$ is distinct: equality of two differences gives equality of two unordered sums, and the Sidon property forces the same ordered endpoints.
For $1\le h<n$, consider differences $a_{i+j}-a_i$ with $1\le j\le h$ and $1\le i\le n-j$. There are
\[
 M=hn-\frac{h(h+1)}2
\]
distinct positive integers, whose sum is at least $M(M+1)/2$. On the other hand, for each $j$,
\[
 \sum_{i=1}^{n-j}(a_{i+j}-a_i)
 =\sum_{i=n-j+1}^n a_i-\sum_{i=1}^j a_i\le jD.
\]
Summing over $j$ proves
\[
 D\ge\frac{M(M+1)}{h(h+1)}.\tag{1}
\]
With $h=\lfloor\sqrt n\rfloor$ for $n\ge3$, this is $D\ge n^2-O(n^{3/2})$, with an absolute constant. Indeed $M=hn(1-O(n^{-1/2}))$ and $h/(h+1)=1-O(n^{-1/2})$. The simpler full-difference count only gives $D\ge\binom n2$.

\section*{Consequences for squared sumset gaps}
There are exactly $t=n(n+1)/2$ sums and their span is $2D$. Cauchy--Schwarz gives
\[
 Q(A)\ge\frac{4D^2}{t(t-1)}.\tag{2}
\]
Inserting (1) yields $Q(A)\ge16-O(n^{-1/2})$ as $n\to\infty$. This is a lower asymptotic estimate, not an assertion of equality or boundedness of $Q$.
It also proves divergence for every sequence with $D/n^2\to\infty$. Conversely, a bounded-$Q$ counterexample sequence must have $D=O(n^2)$, so only sets in a quadratic-length interval need further consideration.

There is a second useful obstruction. For $n\ge3$, the first gap of $A+A$ is $a_2-a_1$, and the last is $a_n-a_{n-1}$. Hence
\[
 Q(A)\ge\frac{(a_2-a_1)^2+(a_n-a_{n-1})^2}{t}.
\]
An endpoint gap growing faster than $n$ therefore also forces divergence. Neither property is proved for every large Sidon set.

\section*{Corrections and remaining gap}
This replaces the attributed GPT-Pro-5.2 draft's unsupported numerical evidence. Its elementary difference-counting expression tends to four, not one; the refined estimate above improves it to sixteen. Its proposed greedy set $\{0,1,3,8,10\}$ is not Sidon, since $1+10=3+8$. Even for $\{0,1\}$ the exact value is $2/3$, so a claimed value $1/3$ cannot be used as a check.
The remaining task is to force unbounded squared-gap mean even when $D=O(n^2)$ and both endpoint gaps are $O(n)$. A single growing gap would be insufficient unless its squared size is large relative to $t\asymp n^2$. All estimates here are elementary reconstructions; no new general solution is claimed.

\textbf{Completion rate estimate: 20\%.}
\end{document}
